\subsection{The perturbed metric and the vertical curves}
\label{subsec:flag-wall-positivity}

Height estimates over number fields control horizontal points. For the arithmetic intersection near a wall we must also control curvature and degrees on the vertical curves.

For the same marked point we compare the smooth reference
height \(h_0(w)\), the interior Hodge height \(h_{\rm H}(w)\), and
the height \(h_M\) of \(M_B(v,s)\) before taking any
intersection. With \(M'\) the marked modulus, the finite set of primes at which the metric changes is
\[
 Q_B=\{p:p\text{ prime},\ p\le B\}
      \mathbin{\cup}\{p:p\text{ prime},\ p\mid 2|\operatorname{disc}E|M'\}.
\] Use the Hodge--Tate lattice at \(p\in Q_B\)
and the de Rham lattice elsewhere.

\paragraph{Smoothing and the finite change of lattice.}
On the complex toroidal compactification use the Kähler reference
form \(\Omega\) and the Hodge potential relative to its smooth
root metric. For a proper initial flag step this is the difference
of the logarithmic determinant metrics. For the relative potential \(u\) and the smooth reference curvature \(\theta\), positivity on the open part reads
\[
 \theta+dd^c u\ge0.
\]
Indeed these fundamental lines are the dual determinants of the first
steps in the Hodge flag, with the conjugate flag in reverse order in
the second block. In period coordinates the curvature on the base
is the Hodge form on a subspace, and variation of the subspace adds
the squared norm of its second fundamental form to the dual
determinant curvature. Past the middle use the conjugate step and
the flat full determinant.

In the toroidal logarithmic charts the
Hodge metric and its inverse have polynomial growth in the boundary
radius
\[
 r=e-\log\|s_D\|^2,
\]
where \(s_D\) cuts out the boundary.
Taking the logarithm of a determinant bounds the potential by
\(C\log r+C\) in absolute value. A locally plurisubharmonic
potential with this sublogarithmic growth extends across a normal
crossing divisor. To see this, add \(\delta\sum\log|z_i|\) at its local branches,
extend across them by the upper limit and take the upper
semicontinuous regularization as \(\delta\downarrow0\).
The resulting potential is bounded above on each compact chart.
For the full central weight use the smooth extension of the
flat determinant potential, whose curvature cancels its reference
curvature identically. Its coefficient may have either sign.
\label{cor:flag-wall-positivity-1}

\begin{proposition}[A fibrewise nef perturbed root metric]
\label{prop:perturbed-flag-metric}
With the permitted smoothing and finite cutoff, \(M_B(v,s)\)
on each of the two selected marked covers has an archimedean semipositive metric. On every vertical curve \(C\) its degree satisfies
\[
 \deg_C(M_B(v,s)|_C)\ge0.
\]
 It satisfies the interior, boundary and dense-set
height comparisons \eqref{eq:flag-actual-ht}.
\end{proposition}

\begin{proof}
\proofstep{Regularizing the fundamental potentials.}
Take \(L_s=s^{-2}\) and smooth the proper fundamental potentials by
the holomorphic part of the Chern exponential map at radius
\(\eta=\exp(-L_s)\). In Demailly's regularization theorem the
background form is closed, its sum with \(dd^c\) of the potential
is a positive current, and the compact Kähler metric admits a
curvature lower bound on orthogonal tangent pairs. A multiple of
\(\Omega\) provides that bound. The regularized potential
\(\psi_\eta=\rho_\eta\psi+K\eta^2\) majorizes \(\psi\) on the
open part. It satisfies, uniformly on the compactification,
\[
 -C(1+\log L_s)\le\psi_\eta\le C,\qquad
 \theta+dd^c\psi_\eta\ge -C\frac{\log L_s}{L_s}\Omega .
\]
Here are the rates in this use of regularization. In each crossing
chart the ball mean of \(-\log |z_i|\) truncated below at zero is
\(O(1+\log(1/\eta))\).

To obtain this bound, rescale the ball when its centre is within
twice the radius of the hyperplane, and use the distance bound
on the other balls. The normalized kernel is bounded by
\(C\eta^{-2d}\) on a ball of radius \(C\eta\) in complex
dimension \(d\). The smooth chart factors change the estimate by
a constant. Jensen's inequality for the logarithm of the boundary
radius gives the lower bound on \(\psi_\eta\). For the curvature
bound Demailly's horizontal Hessian estimate costs
\[
 O(\eta^2+\partial_{\log\eta}\psi_\eta)\Omega.
\]
The radial
Hessian bound gives convexity in \(\log\eta\) after a further
fixed multiple of \(\eta^2\) is added. Comparing with a fixed
small radius gives a secant bound \(O((1+\log L_s)/L_s)\) for
the derivative, and gives the displayed estimate for small \(s\).
There are finitely many fundamental steps and charts for the fixed
model. \cite{demailly-flag}\label{cor:flag-wall-positivity-2}

Form the weighted sum of these potentials by the successive gaps
of \(w\), and its last coefficient for the full determinant. Call
the result \(\phi_s(w)\). In each block of consecutive proper steps the coefficients sum by telescoping, 
\[
 \sum_{i=a}^{b-1}(w_i-w_{i+1})=w_a-w_b\le2\|w\|_\infty.
\]
 They are nonnegative by dominance, so their total on the fixed model is \(O(\|w\|_\infty)\).
The smooth auxiliary line \(A\) has curvature at least \(c\Omega\),
with \(c>0\) on every component. For \(\|w\|_\infty\le M\), choose
\(s_0>0\) so that \(CMs\log(1/s)\le c\) for \(0<s\le s_0\).
The curvature of the perturbed metric then satisfies
\[
 \operatorname{curv} M_B(v,s)
 \ge \bigl(sc-C\|w\|_\infty s^2\log(1/s)\bigr)\Omega
 \ge s\bigl(c-CMs\log(1/s)\bigr)\Omega\ge0 .
\] Its potential is at least the open Hodge
potential on the interior, globally bounded above by a constant,
and globally bounded below by \(-C\log(1/s)-C\).

\proofstep{Finite changes of the root metrics.}
For each root, at a changed prime \(p\) the normalized point-height
difference is at most
\[
   J_p=\frac{\log p}{p-1}+
     \frac{v_p(|\operatorname{disc}E|)}{[E:\Q]}\log p .
\]
Indeed the Hodge--Tate and de Rham integral steps differ by the
Hodge--Tate lattice loss \(1/(p-1)\) and the normalized different
in the labelled projection. The same estimate holds on the reversed
dual roots. In summing local readings on the field of a point,
the degrees above \(p\) add to the degree of that field. The second
sum in \(J_p\) is bounded by \(\log\operatorname{rd}_E\), and
\[
 \sum_{p\le B}\log p/(p-1)=O(1+\log B)
\]
by the elementary
Chebyshev bound and partial summation. The extra marked primes have
fixed total cost. Hence switching the finite lattices from the
empty set to \(Q_B\) costs \(O(\|w\|_\infty\log B)\), uniformly
over the points. \cite{scholze-weinstein-flag}\label{cor:flag-wall-positivity-3}

\proofstep{Assembling the perturbed point height.}
At a point the arithmetic height is additive. Its terms are the
smooth height at \(Q_B\), the average of \(\phi_s(w)\) over the
complex readings of the point, and \(s h_A\). The preceding
comparisons give
\begin{align*}
  h_M&\ge h_{\rm H}(w)+s h_A-C\log B &&\text{in the interior},\\
  h_M&\ge h_0(w)+s h_A-C\log(B/s) &&\text{everywhere},\\
  h_M&\le h_0(w)+s h_A+C\log B &&\text{everywhere}.
\end{align*}
All three use the same finite lattices. Given \(\|v\|_\infty\le1\),
let the middle gaps of \(v\) be at least \(t\), and the other gaps
at least a fixed \(c_0>0\).

Take \(s=t^2\),
\(t\le\min(c_0,1)\) and \(4t\|d_0\|_\infty\le1\). Every
shifted adjacent gap is at least \(t/2\), since it loses at most
\(2t^2\|d_0\|_\infty\). Also \(\|w\|_\infty\le2\).
Use \(b=t/2\) in Section~\ref{subsec:flag-point-heights}. With
\(h=\max(h_F,0)\) the resulting bounds, after enlarging \(C\), are
\begin{align}
 h_M &\ge \tfrac t4 h-C\log(2+h)-C\log B-C
           &&\text{in the interior},\notag\\
 h_M &\ge t^2 h_A-C\log(2+C+h_A)-C\log(B/t^2)-C
           &&\text{on the boundary},\label{eq:flag-actual-ht}\\
 h_M &\le C\log B+C &&\text{on a dense set in each component}.
           \notag
\end{align}
The lower bound for \(h_A\) applies also in the interior.

\paragraph{Degrees at a Hodge--Tate prime.}
For \(p\in Q_B\) pull the root lines to a common normal projective
cover splitting level \(p^2\) and the roots of unity needed for its
torsion characters. Consider a first step of length \(k\) in an
eigenflag. In each affine tangent frame the characters on the
\(p^2\)-torsion have integral representatives for their logarithmic
differentials. Evaluate their \(k\)-fold minors on the step
determinant. At an integral test point the largest minor gives
the Hodge--Tate determinant norm. If \(m_\alpha\) are the normalized minors in the induced frame, integrality and the unit minor give
\[
 \min_\alpha v(m_\alpha)=0.
\]

On overlaps
the differences of the two representatives for each minor have
normalized absolute value strictly less than one. Thus their
residue sections agree. More explicitly, the integral point
criterion on the normal model gives regularity at the special
fibre from the generic regularity and the integral readings.
At a point of that fibre take an integral test point whose
closed specialization detects units in the local ring. A unit
minor there gives a generator on a neighbourhood, and the strict
overlap inequality gives equality of the restrictions to every
reduced vertical curve. The residue sections generate the step line on each such
curve. For a curve \(C'\) above \(C\), let \(L_k\) be the
step line and let \(A_{\rm base}\) account for the discriminant
correction after clearing by \(d_E=[E:\Q]\). The line measured
by the generated minors is
\[
 L_k^{\otimes d_E}|_{C'}
\]

times the pullback of \(A_{\rm base}\). Global generation on \(C'\) gives nonnegative degree, and the
base line pulls back with degree zero. Consequently
\[
 0\le\deg_{C'}(L_k^{\otimes d_E}\otimes A_{\rm base}|_{C'})
    =d_E\deg_{C'}L_k .
\] To descend choose a closed
point of the fibre over the generic point of \(C\) and take
its closure \(C'\); the map has a positive finite degree
\(n\). The projection formula
\[
 \deg_{C'}L_k=n\deg_C L_k
\]

gives the same inequality on \(C\). This construction uses the determinant on the first \(k\)
steps, including at a prime of the discriminant.
\label{cor:flag-wall-positivity-4}

The full signed determinant at each embedding has degree zero on
these curves. The flat determinant section has a constant normalized
Hodge--Tate valuation on the component at the chosen place. The
section and its inverse therefore induce nef lines there. The
comparison uses the determinant frame at that place and its
valuation at a test point on the same component. For a step past
the split \(j\), the negative prefix is a negative prefix of the
conjugate Lie flag plus a full signed determinant. Its degree is
again nonnegative. Clear the rational weight denominators.
For an antitone weight \(w_1\ge\cdots\ge w_N\),
\[
 w=\sum_{k=1}^{N-1}(w_k-w_{k+1})\mathbf1_{[1,k]}
                         +w_N\mathbf1_{[1,N]}.
\]
The gaps are nonnegative, and the full signed determinant
has degree zero on the tested curves, so the root combination
has nonnegative degree. Adding \(sA\) preserves it. \label{cor:flag-wall-positivity-5}

\paragraph{Degrees at a de Rham prime.}
For \(p\notin Q_B\) we have \(p>2\), \(p>B\), and \(p\) divides
neither the field discriminant nor the marked level. Test on an
integral vertical curve. Under a dominating projective model a
contracted curve has zero pullback degree. In the other case
normalize a dominating curve over \(\overline{\mathbb F}_p\).
This multiplies all degrees by the same positive factor. The
unramified action splits into eigenbundles.

In using the semiabelian
family on this curve, one can pass to a finite cover, separate the
constant split torus and take the polarized abelian generic part.
For a subbundle \(W\), let \(W'\) denote its image in the quotient and \(K\) its kernel in the constant torus part. Then
\[
 \deg W=\deg W'+\deg K\le\deg W',
\]
since \(K\) is a subsheaf of a trivial bundle.
The full eigenbundle degrees are preserved under removal of the
constant torus. Further finite pullback multiplies degrees by
the map degree. Thus the flag prefix upper bounds below descend
with the same factor as the total degrees.
\label{cor:flag-wall-positivity-6}

\proofstep{The Hodge and conjugate filtrations.}
Write \(U=\bigoplus_\sigma U_\sigma\) for the Lie bundle of the
resulting family and \(D_\sigma\) for its de Rham eigenbundles.
The polarized Hodge and conjugate filtrations on the abelian
open give
\[
 \begin{aligned}
 0&\to U_{\bar\sigma}^\vee\to D_\sigma\to U_\sigma\to0,\\
 0&\to F^*U_\sigma\to D_{\varphi\sigma}
          \to F^*(U_{\bar\sigma}^\vee)\to0 ,
 \end{aligned}
\]
where \(F\) is Frobenius and \(\varphi\) is its permutation
of the embeddings. For precision, the sequence on the good open must extend as an
exact sequence of vector bundles at each missing point. In the
curve comparison one first makes the finite Frobenius pullback,
using the Frobenius twist of the base field structure on the smooth
curve. Over the abelian open take the two filtrations on primitive
de Rham cohomology with its polarization pairing.

At a missing
point, in the completed discrete valuation ring let \(F_0\) denote
the rational Hodge subspace and \(\Omega_0\) the lattice of integral
invariant differentials. The integral comparison for the
semiabelian torsion tower gives a finitely generated lattice
\(\Lambda\) stable under the endomorphisms with Rosati adjoints.
Its two conditions of adaptation for a rational subspace \(V\)
and lattice \(N\subset V\) are
\[
 \Lambda\cap V=N,\qquad \Lambda+V=N^{\vee,\beta},
\]
where the right side means the inverse image of the integral
linear dual under the polarization pairing \(\beta\).
The integral polarized comparison for the \(p\)-divisible tower
gives both identities at \((F_0,\Omega_0)\). It also gives
\[
 \Lambda\cap W=R\langle\operatorname{Fr}(\Lambda)\rangle,
 \qquad W=K\langle\operatorname{Fr}(H)\rangle ,
\]
where \(R\) is the completed local ring and \(H\) the rational
primitive de Rham space over its fraction field \(K\). This is the
integral Frobenius saturation of the transported lattice.

Since
Frobenius kills \(F_0\) and
\(\Lambda+F_0=\Omega_0^{\vee,\beta}\), this intersection is exactly
\(N=R\langle\operatorname{Fr}(\Omega_0^{\vee,\beta})\rangle\).
The same comparison sends \(\Lambda\) under Verschiebung into the
integral span of the images of \(\Omega_0\) in the Frobenius-twisted
space. In the Frobenius--Verschiebung adjunction, writing that
image of \(x\in\Lambda\) as \(\sum_i a_i\iota(\omega_i)\) gives
\[
     \beta(x,\operatorname{Fr} y)
          =\sum_i a_i\beta(\omega_i,y)^p\in R
                    \qquad (y\in\Lambda).
\]
Here \(a_i\in R\), \(\omega_i\in\Omega_0\), and integrality of the
pairings on the right follows from Hodge adaptation. Thus the
projection of \(\Lambda\) to \(H/W\) lies in the full dual of
\(N\).

Let \(P_W\) be the image of \(\Lambda\) in \(H/W\), and
use the pairing to view it in \(N^\vee\). The exact sequence
\(0\to N\to\Lambda\to P_W\to0\) gives, in a symplectic basis,
\[
 v(\det\Lambda)=v(\det N)+v(\det P_W)
       =\operatorname{length}_R(N^\vee/P_W).
\]
Indeed
\[
 v(\det N^\vee)=-v(\det N)
\]
in the dual basis.
Computing the same determinant with the Lagrangian \(F_0\)
instead gives zero, since Hodge adaptation identifies the
projection with \(\Omega_0^\vee\). The resulting inclusion has zero colength, so
\[
 \operatorname{length}(N^\vee/P_W)=0\quad\Longrightarrow\quad P_W=N^\vee.
\]
This is the second adaptation identity
at \((W,N)\).

The comparisons use the Weil pairing and its extension on the
torsion tower, with functoriality for the paired endomorphisms.
The two adaptation identities give the sublattice and full
quotient in each sequence. To descend a
completed lattice, approximate a basis by rational vectors to
precision strictly finer than its elementary divisors. The change
of basis then belongs to the integral general linear group of
the completion. The rational lattice with this completed span
has the same intersections and quotients by faithful flatness.
It modifies the bundle at just the finite set of missing points.
The two sequences extend there by the prescribed lattices and
the split integral action gives their eigensequences.
\label{cor:flag-wall-positivity-7}
Writing \(d_\sigma=\deg D_\sigma\) gives
\[
 d_{\varphi\sigma}=p d_\sigma.
\]
For an orbit of length \(h\), iterating gives
\[
 d_\sigma=d_{\varphi^h\sigma}=p^h d_\sigma .
\]
Since \(p^h>1\), \(d_\sigma=0\), hence
\(\deg U_\sigma=\deg U_{\bar\sigma}\).

\proofstep{Uniform control by the largest subsheaf degree.}
Let \(P\) be the largest degree of a coherent subsheaf of \(U\).
The degrees are integers. For an upper bound, twist the
dual sufficiently to generate it, embed \(U\) into a sum of fixed
ample lines, and take exterior powers. The zero subsheaf gives
\(P\ge0\). Saturate a subsheaf attaining \(P\). Its degree remains
\(P\). Its Frobenius pullback of degree \(pP\) injects in the full
de Rham bundle \(D\), an extension of \(U\) by \(U^\vee\).
If \(P'\subset U\) realizes \(P\), let \(B\) and \(A\) be
the image and kernel of \(F^*P'\to D\to U\), respectively.
Then \[
 0\to A\to F^*P'\to B\to0,\qquad \deg B\le P.
\]
Saturate \(A\subset U^\vee\). The dual of its locally free
quotient embeds in \(U\), so
\[
 \deg A\le\deg U^\vee-\deg(U^\vee/A^{\rm sat})
                  \le-\deg U+P .
\]

Now \(pP=\deg A+\deg B\le-\deg U+2P\), giving
\[
        (p-2)P\le e,\qquad e=-\deg U .
\]
For the root degrees \(r_{\tau i}\) on the resulting curve, set
\(S_{\tau k}=\sum_{i<k}r_{\tau i}\).
The full sum is the difference of the two conjugate Lie
degrees, so \(S_{\tau N}=0\). The middle sum is one full Lie
eigendegree, hence \(-2\sum_\tau S_{\tau j}=e\). Before the
middle, a prefix is the degree of a flag subsheaf and is at
most \(P\). After the middle, reverse the conjugate flag and
use the vanishing full sum to obtain the same upper bound.
The argument gives these properties for a common positive
multiple of the original readings. Rescale \(P,e\) by the
same factor on descent. \label{cor:flag-wall-positivity-8}

Put \(S=\|w\|_\infty\) and suppose
\[
 4f S\le(B-2)\delta_{\tau j}
\]
for each \(\tau\).
Let \(b_p=4fS/(p-2)\), so that \(\delta_{\tau j}\ge b_p\).
Summation by parts and \(P-S_{\tau k}\ge0\) now give
\[
 \begin{split}
 -\sum_{\tau,i} w_{\tau i}r_{\tau i}
 &=\sum_{\tau,k<N,\ k\ge1}\delta_{\tau k}(P-S_{\tau k})
              -P\sum_\tau(w_{\tau0}-w_{\tau,N-1})\\
 &\ge b_p\bigl(|\mathcal T|P+e/2\bigr)-2fSP\ \ge 0 .
 \end{split}
\]
This proves nonnegative degree at every remaining prime.
After adding \(sA\), clearing the line denominators and descending
from the common model we obtain fibrewise nefness of \(M_B(v,s)\)
on each marked component. Together with the archimedean
semipositivity and \eqref{eq:flag-actual-ht} this supplies the
point and positivity estimates for the arithmetic intersection
on both selected covers. \label{cor:flag-wall-positivity-9}
\end{proof}
